% GENERATED FILE. Edit canonical lessons, then run npm run generate:books.
% canonical-source-hash: fnv1a64:7520b462f31cab77
% canonical-lessons: AR-W07-hook-family-ha-kha, AR-W08-kaf-and-ra, AR-W09-khayr-bikhayr, AR-C03-kayfa, AR-C03-hal, AR-C03-kayfa-haluka, AR-C03-bi-khayr, AR-C03-al-hamdu-lillah, AR-C03-practice

\chapter{How Are You?}
\label{ch:ar-how-are-you}

This chapter is generated from the canonical micro-lessons used by Language Ladder.
Each section stays independently resumable and preserves the authored prerequisite order.

\section[ḥāʾ, khāʾ, jīm]{\ar{ح،} \ar{خ،} \ar{ج} — the dots trick, on a brand-new body}

\label{lesson:AR-W07-hook-family-ha-kha}



\begin{quote}
\textbf{Warm-up.} {[}PAUSE 2s{]} Back in Lesson 4 you learned a whole family of letters that share one \textbf{bowl} and differ only by dots. Arabic does the \emph{same trick} on a \textbf{second} skeleton — a curvy \textbf{hook-and-tail body} — giving three more common letters. Learn the body once; the dots do the rest, exactly as before.
\end{quote}

\subsection*{Script — The hook-and-tail family}
Every letter here is the \textbf{same shape}: a rounded \textbf{hook} at the top flowing into a \textbf{tail that drops below the line}. Only the dot changes:

\noindent
\begin{tabularx}{\linewidth}{@{}>{\raggedright\arraybackslash}X>{\raggedright\arraybackslash}X>{\raggedright\arraybackslash}X@{}}
\toprule
\textbf{letter} & \textbf{sound} & \textbf{the dot} \\
\midrule
\textbf{\ar{ح}} \emph{ḥāʾ} & ḥ — a deep, breathy "h" (from the throat) & \textbf{no} dot \\
\textbf{\ar{خ}} \emph{khāʾ} & kh — like the \emph{ch} in Scottish "lo\textbf{ch}" & \textbf{one} dot \textbf{above} \\
\textbf{\ar{ج}} \emph{jīm} & j — as in "\textbf{j}am" & \textbf{one} dot \textbf{below} \\
\bottomrule
\end{tabularx}

It is the bowl-family idea reborn: one skeleton, the dots deciding everything. Say the dial — \textbf{no dot / above / below $\to$ ḥ / kh / j}.

\subsection*{Script — Draw them}
\begin{itemize}
\raggedright
  \item \textbf{\ar{ح}} (\emph{ḥāʾ}) — the hook body, then the descending tail; \textbf{no dot}. From Phoenician \textbf{ḥēth} ("fence") — the ancestor of Latin \textbf{H}.
  \item \textbf{\ar{خ}} (\emph{khāʾ}) — the \textbf{same} body, plus \textbf{one dot above}. (Arabic built \emph{khāʾ} from \emph{ḥāʾ} by adding the dot — no separate Latin cousin, but you already know the shape.)
  \item \textbf{\ar{ج}} (\emph{jīm}) — the same body, plus \textbf{one dot below}. From Phoenician \textbf{gīml} ("camel") — the great-grandparent of Greek \emph{gamma} and Latin \textbf{C/G}.
\end{itemize}

\subsection*{Script — Where you've seen it — \ar{خير}}
You have already read \textbf{\ar{خير}} (\emph{khayr}, "good/well") inside the Chapter 1 greeting \textbf{\ar{صباح} \ar{الخير}} (\emph{ṣabāḥ al-khayr}, "good morning"). It \textbf{opens with \ar{خ}} — the hook-body plus one dot above. You now know the first letter of a word you've been saying since Lesson… one.

\subsection*{Guided Practice}
{[}PAUSE 1s{]}

\begin{itemize}
\raggedright
  \item {[}YOU WRITE: the same hook-and-tail body three times, then dot them — \ar{ح} (none), \ar{خ} (above), \ar{ج} (below){]}
  \item {[}YOU SAY: the dial — "no dot, above, below — ḥ, kh, j"{]}
  \item {[}YOU SAY: "\ar{خ} opens \ar{خير} — the 'good' in \emph{ṣabāḥ al-khayr}"{]}
\end{itemize}

\begin{tcolorbox}[breakable,colback=teal!4,colframe=teal!35!black,title={Wrap-up recall}]
{[}PAUSE 3s{]} What do \textbf{\ar{ح}, \ar{خ}, \ar{ج}} share, and what tells them apart? (One \textbf{hook-and-tail body}; the \textbf{dots} — none / above / below.) Which sound is \emph{\ar{خ}}, and where's its dot? (\textbf{kh}, like \emph{loch}; \textbf{one above}.) Which Chapter 1 word opens with \emph{\ar{خ}}? (\textbf{\ar{خير}}, in \emph{ṣabāḥ al-khayr}.) Next: \textbf{\ar{ك}} (\emph{kāf}) and \textbf{\ar{ر}} (\emph{rā}) — the last letters you need to write it.
\end{tcolorbox}
\section[kāf, rā]{\ar{ك} and \ar{ر} — an angular k, and an r that won't hold hands}

\label{lesson:AR-W08-kaf-and-ra}



\begin{quote}
\textbf{Warm-up.} {[}PAUSE 2s{]} Two more shapes, and then — next lesson — you can write \textbf{\ar{خير}}. Meet \textbf{\ar{ك}} (\emph{kāf}, "k") and \textbf{\ar{ر}} (\emph{rā}, "r"). One is angular and new; the other revisits a rule you learned on the very first day.
\end{quote}

\subsection*{Script — \ar{ك} (kāf) — the angular k}
Break it apart:

1. \textbf{the body} — an \textbf{angular upright}, bending like an open corner (not a curve — a hard angle). 2. \textbf{the inner stroke} — a small \emph{hamza}-like flick tucked \textbf{inside} the angle.

\textbf{\ar{ك}} is \emph{kāf}, a clean "k." From Phoenician \textbf{kaph} ("palm of the hand") — the ancestor of Greek \emph{kappa} and Latin \textbf{K}. It joins its neighbours, changing shape by position like the other connectors.

\subsection*{Script — \ar{ر} (rā) — the little downward curve that breaks}
Break it apart — there's just \emph{one piece}:

1. \textbf{the curve} — a short stroke that \textbf{curves down below the baseline} and stops. No dot, no loop.

\textbf{\ar{ر}} is \emph{rā}, "r." From Phoenician \textbf{rēsh} ("head") — ancestor of Greek \emph{rho} and Latin \textbf{R}. And here is the important rule: like \textbf{alif} (Lesson 1), \textbf{rā does not join to the letter after it} — it is one of the \textbf{six non-joiners}. So \emph{rā} forces the pen to \textbf{lift}, just as \emph{alif} did in \emph{\ar{سلام}}.

\subsection*{Guided Practice}
{[}PAUSE 1s{]}

\begin{itemize}
\raggedright
  \item {[}YOU WRITE: \textbf{\ar{ك}} — angular body + the inner flick{]}
  \item {[}YOU WRITE: \textbf{\ar{ر}} — one short downward curve below the line{]}
  \item {[}YOU SAY: "\emph{rā}, like \emph{alif}, \textbf{breaks} — no joining the next letter"{]}
  \item {[}YOU SAY: "\emph{kaph} $\to$ K, \emph{rēsh} $\to$ R — the Phoenician cousins again"{]}
\end{itemize}

\begin{tcolorbox}[breakable,colback=teal!4,colframe=teal!35!black,title={Wrap-up recall}]
{[}PAUSE 3s{]} Describe \textbf{\ar{ك}}: angular or curved, and what sits inside? (\textbf{Angular} body with a small \textbf{inner stroke}.) What Latin letters descend from \emph{kāf} and \emph{rā}? (\textbf{K} and \textbf{R}.) What rule does \textbf{\ar{ر}} share with \textbf{alif}? (It is a \textbf{non-joiner} — the pen \textbf{lifts} after it.) Next: you assemble \textbf{\ar{خير}} — and answer "how are you?"
\end{tcolorbox}
\section[khayr, bi-khayr]{\ar{خير} and \ar{بخير} — writing "good," and how you say you're well}

\label{lesson:AR-W09-khayr-bikhayr}



\begin{quote}
\textbf{Warm-up.} {[}PAUSE 2s{]} Everything comes together. With \textbf{\ar{خ}} (Lesson 7), \textbf{\ar{ي}} (Lesson 5), and \textbf{\ar{ر}} (Lesson 8), you can now write \textbf{\ar{خير}} (\emph{khayr}, "good") — the word hiding in every "good morning" you've greeted with — and then \textbf{\ar{بخير}} (\emph{bi-khayr}, "well"), the everyday answer to "how are you?"
\end{quote}

\subsection*{Script — Assemble \ar{خير} — right to left}
Write it starting on the \textbf{right}, moving \textbf{left}:

1. \textbf{\ar{خ}} (\emph{khāʾ}) — the hook-body with \textbf{one dot above}, \textbf{reaching left} into… 2. \textbf{\ar{ي}} (\emph{yāʾ}) — joined in the middle (its two dots below), flowing into… 3. \textbf{\ar{ر}} (\emph{rā}) — the little down-curve. Remember Lesson 8: \emph{rā} is a \textbf{non-joiner}, so it simply ends the word — but \emph{yāʾ} \textbf{does} join into it from the right.

Read back, right to left: \textbf{\ar{خ} · \ar{ي} · \ar{ر}} = \emph{kh · ī · r} = \textbf{\ar{خير}}, "good." (The root \textbf{kh–y–r} means "good, best, choice" across Semitic — the same three consonants carry the meaning, vowels laid on top, just like Arabic always works.)

\subsection*{Script — Then \ar{بخير} — "(I'm) well"}
Glue \textbf{\ar{بـ}} (\emph{bi-}, "in / with") onto the front — \emph{bāʾ} from Lesson 3, its one dot below — and it joins straight into the \emph{khāʾ}:

\begin{itemize}
\raggedright
  \item \textbf{\ar{بخير}} = \emph{bāʾ · khāʾ · yāʾ · rā} = \textbf{bi-khayr}, literally "\textbf{in goodness}" — what you say when someone asks \emph{kayfa ḥāluk?} ("how are you?"). "\emph{Al-ḥamdu lillāh, bi-khayr}" — "Praise God, (I'm) well."
\end{itemize}

You have now written a \textbf{real reply} in a real conversation, in your own hand.

\subsection*{Guided Practice}
{[}PAUSE 1s{]}

\begin{itemize}
\raggedright
  \item {[}YOU WRITE: \textbf{\ar{خير}}, right to left — "khāʾ (dot above), yāʾ, rā — lift"{]}
  \item {[}YOU WRITE: \textbf{\ar{بخير}} — add \emph{bi-} on the right: "bāʾ joins into khāʾ"{]}
  \item {[}YOU SAY: "the root \textbf{kh-y-r} = good; \emph{bi-khayr} = 'in goodness' = 'I'm well'"{]}
\end{itemize}

\begin{tcolorbox}[breakable,colback=teal!4,colframe=teal!35!black,title={Wrap-up recall}]
{[}PAUSE 3s{]} Spell \textbf{\ar{خير}} letter by letter, right to left. (\emph{khāʾ · yāʾ · rā} — "good.") Which letter ends it, and does anything join \emph{after} it? (\textbf{rā} — \textbf{no}, it's a non-joiner; the pen lifts.) What does \textbf{\ar{بخير}} literally mean, and when do you say it? ("\textbf{In goodness}" — the answer to "how are you?", \emph{kayfa ḥāluk?}.) You can now hand-write a greeting \emph{and} its reply — the Arabic script is yours.
\end{tcolorbox}
\section[kayfa]{\ar{كيف} (kayfa) — "how"}

\label{lesson:AR-C03-kayfa}



\begin{quote}
\textbf{Warm-up.} {[}PAUSE 2s{]} Chapter 2 gave you \textbf{\ar{ما}} (\emph{mā}, "what"). Here is its partner: \textbf{\ar{كيف}} (\emph{kayfa}, "how"). Two question words, and you can already build with both.
\end{quote}

\subsection*{You'll want to know — The letters in this word}
\emph{(Skim if you read Arabic.)} \textbf{\ar{ك}} (\emph{kāf}) and \textbf{\ar{ي}} (\emph{yāʾ}) you have written — \emph{kāf} in the writing set, \emph{yāʾ} in \textbf{\ar{اسمي}}. The last letter, \textbf{\ar{ف}} (\emph{fāʾ}, the \emph{f}), is \textbf{new}; the writing track hasn't reached it yet, so read it for now and draw it later.

Right to left: \textbf{\ar{ك} · \ar{ي} · \ar{ف}} $\to$ \emph{kay-fa}.

Note the \textbf{ay} in the middle — the \emph{yāʾ} is doing double duty again, this time as part of a diphthong, exactly as it did as a long \emph{ī} in \emph{ismī}.

\begin{cousinweb}
\textbf{\ar{كيف}} is a \textbf{question word}, so it doesn't behave like the root-and-pattern nouns you've been meeting — it just \emph{is}. But the letters \textbf{k–y–f} do form a root elsewhere:

\noindent
\begin{tabularx}{\linewidth}{@{}>{\raggedright\arraybackslash}X>{\raggedright\arraybackslash}X@{}}
\toprule
\textbf{word} & \textbf{meaning} \\
\midrule
\textbf{kayfiyya} & quality, manner — literally "\textbf{how-ness}" \\
\textbf{kayyafa} & to adapt, adjust (to make a thing fit \emph{how} it must be) \\
\textbf{takyīf} & adaptation — and, in modern Arabic, \textbf{air conditioning} \\
\bottomrule
\end{tabularx}

So the abstract noun for "quality" is built straight from the word "how." Arabic does this constantly: take a word, run it through a pattern, get a concept.
\end{cousinweb}

\subsection*{Guided Practice}
{[}PAUSE 1s{]}

\begin{itemize}
\raggedright
  \item {[}YOU SAY: "\emph{kayfa}" — "how"{]}
  \item {[}YOU SAY: the pair — "\textbf{\ar{مَا}} \emph{mā} (what) · \textbf{\ar{كيف}} \emph{kayfa} (how)"{]}
  \item {[}YOU SAY: "\emph{kayfiyya}" — and hear "how-ness"{]}
\end{itemize}

\begin{tcolorbox}[breakable,colback=teal!4,colframe=teal!35!black,title={Wrap-up recall}]
{[}PAUSE 3s{]} What does \textbf{\ar{كيف}} mean? ("\textbf{How}.") Which Chapter 2 word is its partner? (\textbf{\ar{ما}}, "what".) Which letter in it is new? (\textbf{\ar{ف}}, \emph{fāʾ} — the \emph{f}.) What is \emph{kayfiyya}? ("\textbf{Quality}" — literally \emph{how-ness}, built from this very word.) Next: the word Arabic asks about when it asks how you are — and it isn't "you."
\end{tcolorbox}
\section[ḥāl]{\ar{حال} (ḥāl) — "state, condition"}

\label{lesson:AR-C03-hal}



\begin{quote}
\textbf{Warm-up.} {[}PAUSE 2s{]} English asks "how are \textbf{you}?" Arabic asks about something else entirely — your \textbf{\ar{حال}}, your \emph{state}. Get this word and the whole question falls open.
\end{quote}

\subsection*{You'll want to know — The letters in this word}
\emph{(Skim if you read Arabic.)} Three letters, \textbf{all already yours}: \textbf{\ar{ح}} (\emph{ḥāʾ}, the deep throat \emph{h}, from the hook-family lesson), \textbf{\ar{ا}} (\emph{alif}), \textbf{\ar{ل}} (\emph{lām}).

Right to left: \textbf{\ar{ح} · \ar{ا} · \ar{ل}} $\to$ \emph{ḥāl}. And remember from the writing set: \emph{alif} \textbf{breaks} the join, so \emph{ḥā-} then \emph{-l}.

\begin{culture}
Here is the payoff. \textbf{\ar{حال}} comes from the root \textbf{ḥ–w–l}, "\textbf{to turn, to change}." Look what grows from it:

\noindent
\begin{tabularx}{\linewidth}{@{}>{\raggedright\arraybackslash}X>{\raggedright\arraybackslash}X>{\raggedright\arraybackslash}X@{}}
\toprule
\textbf{word} & \textbf{meaning} & \textbf{the turning inside} \\
\midrule
\textbf{ḥāl} & state, condition & how things have \textbf{turned out} \\
\textbf{ḥawla} & around & turning \emph{round} something \\
\textbf{taḥwīl} & transformation, transfer & making something turn into another \\
\textbf{ḥawl} & a year & one full \textbf{turn} of the seasons \\
\textbf{aḥwāl} & circumstances (plural of \emph{ḥāl}) & the way things stand \\
\bottomrule
\end{tabularx}

So a "state" in Arabic is not a fixed thing you \emph{are} — it's the position you have \textbf{turned to}. The language builds "condition" out of motion.

Compare what you already know: Spanish \emph{estar} ("to be, temporarily") came from Latin \emph{stāre}, "to \textbf{stand}." Two languages, two metaphors for the same idea — one \textbf{stands}, one \textbf{turns}.
\end{culture}

\subsection*{Guided Practice}
{[}PAUSE 1s{]}

\begin{itemize}
\raggedright
  \item {[}YOU SAY: "\emph{ḥāl}" — deep \emph{ḥ} from the throat, not an English \emph{h}{]}
  \item {[}YOU SAY: "\emph{ḥāl · ḥawla · taḥwīl}" — hear the root turning{]}
  \item {[}YOU SAY: the contrast — "Arabic \textbf{turns}, Spanish \emph{estar} \textbf{stands}"{]}
\end{itemize}

\begin{tcolorbox}[breakable,colback=teal!4,colframe=teal!35!black,title={Wrap-up recall}]
{[}PAUSE 3s{]} What does \textbf{\ar{حال}} mean? ("\textbf{State, condition}.") What root is it from, and what does that root mean? (\textbf{ḥ–w–l}, "\textbf{to turn, change}.") So what is a "state," literally? (\textbf{How things have turned}.) What is \emph{ḥawl}? ("\textbf{A year}" — one full turn.) Next: bolt a "your" onto it and you have the question.
\end{tcolorbox}
\section[kayfa ḥāluka / kayfa ḥāluki]{\ar{كيف} \ar{حالك؟} (kayfa ḥāluka?) — "how are you?"}

\label{lesson:AR-C03-kayfa-haluka}



\begin{quote}
\textbf{Warm-up.} {[}PAUSE 2s{]} Everything in this question is already yours. You are not learning a new phrase — you are \textbf{bolting three known pieces together}.
\end{quote}

\subsection*{You'll want to know — The assembly}
\noindent
\begin{tabularx}{\linewidth}{@{}>{\raggedright\arraybackslash}X>{\raggedright\arraybackslash}X@{}}
\toprule
\textbf{piece} & \textbf{where you met it} \\
\midrule
\textbf{\ar{كيف}} \emph{kayfa} — "how" & last-but-one lesson \\
\textbf{\ar{حال}} \emph{ḥāl} — "state" & last lesson \\
\textbf{\ar{ـك}} \emph{-ka / -ki} — "your" & \textbf{new}, but see below \\
\bottomrule
\end{tabularx}

$\to$ \textbf{\ar{كيف} \ar{حالك؟}} = "\textbf{how} {[}is{]} \textbf{your state}?"

\begin{grammarlens}[title={Grammar Lens — The suffix you already half-know}]
In Chapter 2 you attached \textbf{\ar{ـي}} (\emph{-ī}, "my") to a noun: \emph{ism} $\to$ \textbf{ismī} ("my name"). The suffix for "\textbf{your}" works \textbf{identically} — it just clips onto the end:

\begin{itemize}
\raggedright
  \item \textbf{\ar{حالك}} \emph{ḥāluka} — "your state," said \textbf{to a man}
  \item \textbf{\ar{حالك}} \emph{ḥāluki} — "your state," said \textbf{to a woman}
\end{itemize}

And look at that split: \textbf{-ka} to a man, \textbf{-ki} to a woman. That is exactly the \textbf{\ar{أنتَ} / \ar{أنتِ}} (\emph{anta / anti}) division from Chapter 2 — Arabic marking gender, not formality, where Spanish and French mark register. The same \emph{fatḥa} / \emph{kasra} mark that flipped \emph{anta} to \emph{anti} flips \emph{ḥāluka} to \emph{ḥāluki}. \textbf{One letter, one mark, both meanings.}
\end{grammarlens}

\begin{grammarlens}[title={Grammar Lens — No verb at all}]
Notice what is \textbf{missing}: there is no "is," no "are." Just \emph{how — your-state}. This is the \textbf{zero copula} you met with \textbf{\ar{اسمي}} ("my name {[}is{]} …"): Arabic simply doesn't need a verb to join a thing to its description.
\end{grammarlens}

\subsection*{Guided Practice}
{[}PAUSE 1s{]}

\begin{itemize}
\raggedright
  \item {[}YOU SAY: to a man — "\emph{kayfa ḥāluka?}"{]}
  \item {[}YOU SAY: to a woman — "\emph{kayfa ḥāluki?}"{]}
  \item {[}YOU SAY: the parallel — "\emph{ism\textbf{ī}} my name · \emph{ḥāl\textbf{uka}} your state"{]}
  \item {[}YOU SAY: the missing word — "how · your-state — \textbf{no 'is'}"{]}
\end{itemize}

\begin{tcolorbox}[breakable,colback=teal!4,colframe=teal!35!black,title={Wrap-up recall}]
{[}PAUSE 3s{]} What does \textbf{\ar{كيف} \ar{حالك؟}} literally say? ("\textbf{How {[}is{]} your state?}") Which suffix means "your," and how does it split? (\textbf{-ka} to a man, \textbf{-ki} to a woman — the \emph{anta/anti} split again.) Which Chapter 2 suffix does it copy? (\textbf{-ī}, "my," as in \emph{ismī}.) What word is absent from the sentence? (\textbf{"Is"} — the zero copula.) Next: how to answer.
\end{tcolorbox}
\section[bi-khayr]{\ar{بخير} (bi-khayr) — "well"}

\label{lesson:AR-C03-bi-khayr}



\begin{quote}
\textbf{Warm-up.} {[}PAUSE 2s{]} The answer to \emph{kayfa ḥāluka?} — and you have \textbf{already written this word by hand}. Nothing here is new; it is two old pieces stuck together.
\end{quote}

\subsection*{You'll want to know — The letters in this word}
\emph{(Skim if you read Arabic.)} \textbf{\ar{ب}} (\emph{bāʾ}), \textbf{\ar{خ}} (\emph{khāʾ}), \textbf{\ar{ي}} (\emph{yāʾ}), \textbf{\ar{ر}} (\emph{rāʾ}) — every one of them from the writing set, where you assembled \textbf{\ar{خير}} and then \textbf{\ar{بخير}} letter by letter. Right to left: \textbf{\ar{ب} · \ar{خ} · \ar{ي} · \ar{ر}}.

\subsection*{You'll want to know — The two pieces}
\textbf{\ar{خير}} (\emph{khayr}) — "\textbf{good, goodness}." You have met it before without being told: it is sitting inside the Chapter 1 greeting \textbf{\ar{صباح} \ar{الخير}} (\emph{ṣabāḥ al-khayr}, "morning of \textbf{the good}"). The greeting was never a fixed lump — it was always \emph{ṣabāḥ} + \emph{al-} + \emph{khayr}.

\textbf{\ar{بـ}} (\emph{bi-}) — "\textbf{in, with}." A \textbf{one-letter preposition}, and like \textbf{\ar{الـ}} (\emph{al-}) from Chapter 1, it does not stand alone: it \textbf{glues} onto the front of the next word.

$\to$ \textbf{\ar{بخير}} = \emph{bi-} + \emph{khayr} = "\textbf{in goodness}."

\subsection*{You'll want to know — What you actually say}
\begin{itemize}
\raggedright
  \item — \textbf{\ar{كيف} \ar{حالك؟}} \emph{kayfa ḥāluka?} ("How is your state?")
  \item — \textbf{\ar{بخير،} \ar{شكراً}.} \emph{bi-khayr, shukran.} ("Well, thank you.")
\end{itemize}

Literally: "\emph{In goodness, thanks.}" Not "I am fine" — Arabic doesn't need the "I am" any more than the question needed an "is."

\subsection*{You'll want to know — The bi- you have heard before}
That same \textbf{\ar{بـ}} opens the most-spoken phrase in the language: \textbf{\ar{بسم} \ar{الله}} (\emph{bismillāh}, "\textbf{in} the name of God") — \emph{bi-} + \emph{ism} (the very word from Chapter 2!) + \emph{allāh}. One tiny letter, doing the same job.

\subsection*{Guided Practice}
{[}PAUSE 1s{]}

\begin{itemize}
\raggedright
  \item {[}YOU SAY: "\emph{bi-khayr}" — "in goodness"{]}
  \item {[}YOU SAY: the callback — "\emph{ṣabāḥ al-\textbf{khayr}} … \emph{bi-\textbf{khayr}}"{]}
  \item {[}YOU SAY: the exchange — "\emph{kayfa ḥāluka?} — \emph{bi-khayr, shukran.}"{]}
  \item {[}YOU SAY: "\emph{bismillāh}" — and find \emph{ism} inside it{]}
\end{itemize}

\begin{tcolorbox}[breakable,colback=teal!4,colframe=teal!35!black,title={Wrap-up recall}]
{[}PAUSE 3s{]} What does \textbf{\ar{بخير}} literally mean? ("\textbf{In goodness}.") Where had you already met \emph{khayr}? (Inside \textbf{\ar{صباح} \ar{الخير}} — "morning of the good.") What kind of word is \textbf{\ar{بـ}}? (A \textbf{one-letter preposition} that attaches, like \emph{al-}.) What is missing from the answer, again? (Any word for "\textbf{I am}.") Next: the answer Arabic gives even more often.
\end{tcolorbox}
\section[al-ḥamdu lillāh]{\ar{الحمد} \ar{لله} (al-ḥamdu lillāh) — "praise be to God"}

\label{lesson:AR-C03-al-hamdu-lillah}



\begin{quote}
\textbf{Warm-up.} {[}PAUSE 2s{]} Ask an Arabic speaker how they are and this, more often than \emph{bi-khayr}, is what comes back. It also hides the best root story in the chapter.
\end{quote}

\subsection*{You'll want to know — The letters}
\emph{(Skim if you read Arabic.)} \textbf{\ar{ا} \ar{ل} \ar{ح} \ar{م} \ar{د}} and \textbf{\ar{ل} \ar{ل} \ar{ه}} — \textbf{every letter is one you have written}: \emph{alif}, \emph{lām}, \emph{ḥāʾ}, \emph{mīm}, \emph{dāl}, \emph{hāʾ}. Nothing new to draw.

\subsection*{You'll want to know — The pieces}
\noindent
\begin{tabularx}{\linewidth}{@{}>{\raggedright\arraybackslash}X>{\raggedright\arraybackslash}X>{\raggedright\arraybackslash}X@{}}
\toprule
\textbf{piece} & \textbf{meaning} & \textbf{where from} \\
\midrule
\textbf{\ar{الـ}} \emph{al-} & "the" & Chapter 1 \\
\textbf{\ar{حمد}} \emph{ḥamd} & praise & \textbf{new}, root \textbf{ḥ–m–d} \\
\textbf{\ar{لـ}} \emph{li-} & "to, for" & a one-letter preposition, like \emph{bi-} \\
\textbf{\ar{الله}} \emph{allāh} & God & \emph{al-} + \emph{ilāh}, "\textbf{the} god" \\
\bottomrule
\end{tabularx}

$\to$ \textbf{\ar{الحمد} \ar{لله}} = "\textbf{the praise} {[}is{]} \textbf{to God}." Zero copula once more.

Note \textbf{\ar{الله}} itself: it is \emph{al-} (the article you learned in Chapter 1) fused onto \emph{ilāh}, "a god." The word is literally "\textbf{the God}." Its Semitic cousins are Hebrew \emph{\textbf{El}} and \emph{\textbf{Elohim}} — the same family pairing as \emph{salām} / \emph{shalom} from Chapter 1.

\begin{culture}
\textbf{ḥ–m–d} means "\textbf{to praise}." Run it through Arabic's patterns:

\noindent
\begin{tabularx}{\linewidth}{@{}>{\raggedright\arraybackslash}X>{\raggedright\arraybackslash}X@{}}
\toprule
\textbf{word} & \textbf{pattern meaning} \\
\midrule
\textbf{ḥamd} & praise (the plain noun) \\
\textbf{maḥmūd} & praised \\
\textbf{ḥamīd} & praiseworthy \\
\textbf{aḥmad} & "more praiseworthy" — the name \textbf{Ahmad} \\
\textbf{muḥammad} & "the \textbf{much-praised} one" — the name \textbf{Muhammad} \\
\bottomrule
\end{tabularx}

So the most frequently used phrase in the language and the most common given name on earth are built from \textbf{the same three consonants}. This is the root engine doing exactly what Chapter 1 promised: three letters, and a whole family of meanings drops out.
\end{culture}

\subsection*{You'll want to know — Using it}
\begin{itemize}
\raggedright
  \item — \textbf{\ar{كيف} \ar{حالك؟}} \emph{kayfa ḥāluka?}
  \item — \textbf{\ar{الحمد} \ar{لله}.} \emph{al-ḥamdu lillāh.}
\end{itemize}

It works as "fine, thanks" — said whether things are going well or badly, which is part of its point.

\subsection*{Guided Practice}
{[}PAUSE 1s{]}

\begin{itemize}
\raggedright
  \item {[}YOU SAY: "\emph{al-ḥamdu lillāh}"{]}
  \item {[}YOU SAY: the root — "\emph{ḥamd · ḥamīd · aḥmad · muḥammad}"{]}
  \item {[}YOU SAY: "\emph{allāh}" — and hear \emph{al-} + \emph{ilāh}, "\textbf{the} god"{]}
  \item {[}YOU SAY: the full exchange — "\emph{kayfa ḥāluka?} — \emph{al-ḥamdu lillāh.}"{]}
\end{itemize}

\begin{tcolorbox}[breakable,colback=teal!4,colframe=teal!35!black,title={Wrap-up recall}]
{[}PAUSE 3s{]} What does \textbf{\ar{الحمد} \ar{لله}} literally say? ("\textbf{The praise {[}is{]} to God}.") What root is \emph{ḥamd} from? (\textbf{ḥ–m–d}, "to praise".) Which two names come out of that root? (\textbf{Aḥmad} and \textbf{Muḥammad} — "the much-praised.") What is \emph{allāh} made of? (\emph{\textbf{al-}} + \emph{\textbf{ilāh}} = "\textbf{the god}"; cousin of Hebrew \emph{El}.) Next: practice the whole exchange.
\end{tcolorbox}
\section[Practice]{Practice — the whole exchange}

\label{lesson:AR-C03-practice}



\begin{quote}
\textbf{Warm-up.} {[}PAUSE 2s{]} Chapter 1 let you greet. Chapter 2 let you introduce yourself. This chapter closes the loop: you can now \textbf{ask after someone and answer back}.
\end{quote}

\subsection*{The exchange}
\noindent
\begin{tabularx}{\linewidth}{@{}>{\raggedright\arraybackslash}X>{\raggedright\arraybackslash}X>{\raggedright\arraybackslash}X@{}}
\toprule
\textbf{} & \textbf{Arabic} & \textbf{literally} \\
\midrule
A & \textbf{\ar{السلام} \ar{عليكم}} & peace {[}be{]} upon you \\
B & \textbf{\ar{وعليكم} \ar{السلام}} & and upon you, peace \\
A & \textbf{\ar{كيف} \ar{حالك؟}} & how {[}is{]} your state? \\
B & \textbf{\ar{الحمد} \ar{لله،} \ar{بخير}.} & the praise {[}is{]} to God — in goodness \\
B & \textbf{\ar{وأنت؟} \ar{ما} \ar{اسمك؟}} & and you? what {[}is{]} your name? \\
\bottomrule
\end{tabularx}

Read down the right-hand column. \textbf{Not one line contains a verb "to be."} Arabic joins things by putting them next to each other.

\subsection*{You'll want to know — the pieces}
\noindent
\begin{tabularx}{\linewidth}{@{}>{\raggedright\arraybackslash}X>{\raggedright\arraybackslash}X@{}}
\toprule
\textbf{phrase} & \textbf{built from} \\
\midrule
\emph{kayfa ḥāluka} & \emph{kayfa} + \emph{ḥāl} + \textbf{-ka/-ki} \\
\emph{bi-khayr} & \textbf{bi-} + \emph{khayr} (already inside \emph{ṣabāḥ al-khayr}) \\
\emph{al-ḥamdu lillāh} & \emph{al-} + \emph{ḥamd} + \textbf{li-} + \emph{allāh} \\
\emph{ismuka} & \emph{ism} + \textbf{-ka/-ki} — the same suffix as \emph{ḥāluka} \\
\bottomrule
\end{tabularx}

Three tiny attaching pieces — \textbf{al-}, \textbf{bi-}, \textbf{li-} — and two suffixes, \textbf{-ī} and \textbf{-ka/-ki}. That is most of the chapter.

\begin{grammarlens}[title={Grammar Lens — gender}]
Every "you" in this chapter splits. Say the pair aloud:

\noindent
\begin{tabularx}{\linewidth}{@{}>{\raggedright\arraybackslash}X>{\raggedright\arraybackslash}X@{}}
\toprule
\textbf{to a man} & \textbf{to a woman} \\
\midrule
\emph{kayfa ḥāl\textbf{uka}?} & \emph{kayfa ḥāl\textbf{uki}?} \\
\emph{mā ism\textbf{uka}?} & \emph{mā ism\textbf{uki}?} \\
\emph{ant\textbf{a}} & \emph{ant\textbf{i}} \\
\bottomrule
\end{tabularx}

One vowel mark decides — the \emph{fatḥa} / \emph{kasra} pair from the writing set.
\end{grammarlens}

\subsection*{Guided Practice}
{[}PAUSE 1s{]}

\begin{itemize}
\raggedright
  \item {[}YOU SAY: the full exchange, both parts, to a man{]}
  \item {[}YOU SAY: it again, to a woman — every \emph{-ka} becomes \emph{-ki}{]}
  \item {[}YOU WRITE: \textbf{\ar{بخير}} — you have every letter{]}
  \item {[}YOU SAY: the roots met so far — "\textbf{s-l-m} peace · \textbf{ḥ-w-l} turn · \textbf{kh-y-r} good · \textbf{ḥ-m-d} praise"{]}
\end{itemize}

\begin{tcolorbox}[breakable,colback=teal!4,colframe=teal!35!black,title={Wrap-up recall}]
{[}PAUSE 3s{]} Ask "how are you?" to a woman. (\emph{Kayfa ḥāluki?}) Give both answers. (\emph{Bi-khayr} / \emph{al-ḥamdu lillāh}.) What do \emph{al-}, \emph{bi-} and \emph{li-} have in common? (All \textbf{attach} to the next word — they never stand alone.) What single word is absent from every line of the dialogue? (\textbf{"Is"} — the zero copula.) Which root gave both a common phrase and a common name? (\textbf{ḥ–m–d} — \emph{al-ḥamdu lillāh} and \emph{Muḥammad}.) Next chapter: \textbf{farewells} — and the root that started it all comes back.
\end{tcolorbox}
